% =============================================================================
% ESEMPIO DI APPENDICE FACOLTATIVA - ITALIANO
% Disabilitarla in configuration.tex con:
% \providecommand{\ISISAppendixChoice}{disabled}
% Sostituire il contenuto dimostrativo con materiale di supporto utile ma non
% essenziale per comprendere il corpo principale della tesi.
% =============================================================================

\chapter{Materiale sperimentale aggiuntivo}\label{app:additional-material}

\begin{keypoint}[Appendice facoltativa]
Questa appendice \`e facoltativa e va mantenuta solo quando aggiunge dettagli
di supporto utili. Pu\`o contenere dimostrazioni complete, tabelle estese,
questionari, note implementative, figure aggiuntive, dizionari dei dati e
istruzioni per la riproduzione degli esperimenti.
\end{keypoint}

\section{Protocollo sperimentale esteso}
Riportare il protocollo completo necessario a ripetere la valutazione. In una
tesi reale, sostituire i valori seguenti con la versione esatta del dataset, la
procedura di suddivisione, il preprocessing, i seed casuali, l'ambiente software
e i criteri di arresto.

\begin{table}[htbp]
\centering
\caption{Esempio di configurazione sperimentale.}
\label{tab:appendix-config}
\begin{tabularx}{\textwidth}{lY}
\toprule
Voce & Valore di esempio\\
\midrule
Versione del dataset & Versione 2.1, scaricata il 15 marzo 2026\\
Suddivisione dei dati & 70\% training, 15\% validation, 15\% test; stratificata per classe\\
Seed casuali & 11, 23, 47, 71 e 101\\
Software & Python 3.12; versioni di framework e librerie salvate in un lock file\\
Hardware & Una GPU con 24 GB di memoria; 16 core CPU; 64 GB di RAM\\
Criterio di arresto & Early stopping sulla loss di validation con patience 12\\
\bottomrule
\end{tabularx}
\end{table}

\section{Esempio di comando di riproduzione}
Un comando compatto \`e utile quando il repository contiene uno script capace
di rigenerare i risultati principali.

\begin{verbatim}
python reproduce.py \
  --config configs/final.yaml \
  --seeds 11 23 47 71 101 \
  --output results/final
\end{verbatim}

\section{Risultati aggiuntivi}
Utilizzare questa sezione per analisi secondarie che supportano la discussione
senza interrompere la narrazione principale. Distinguere chiaramente le analisi
confermative da quelle esplorative e motivare eventuali deviazioni dal
protocollo originario.

\begin{resultbox}[Esempio di interpretazione]
Nelle cinque esecuzioni illustrative, l'ordinamento principale dei metodi
confrontati rimane invariato. Sostituire questo testo con la stima reale,
l'intervallo di incertezza e il riferimento alla tabella o figura pertinente.
\end{resultbox}
